\documentclass{beamer}
\usefonttheme{professionalfonts}
%\usetheme{Rochester}
%\usetheme{AnnArbor}
\usetheme{Szeged}
\useoutertheme{infolines}
\usecolortheme{dolphin}
\usepackage{graphicx}
% \usepackage[english]{babel}
\usepackage{multicol}
\usepackage{color}
\usepackage{xcolor}
%\usepackage{smartdiagram}


%\definecolor{carmine}{rgb}{0.59, 0.0, 0.09}
\usepackage{amsmath} 

\usepackage{graphicx}
\usepackage{tikz}
\usetikzlibrary{shapes.callouts}
\usepackage{multirow}

\usepackage[utf8]{inputenc}

\usepackage{enumerate}
\usepackage[shortlabels]{enumitem}
\usepackage{amsmath,amsthm,amssymb,amsfonts,mathrsfs,latexsym,stmaryrd}
\usepackage{tcolorbox}

\usepackage{color} 
\usepackage{amssymb, amsmath, amsbsy}
\usepackage{adjustbox,lipsum}
\usepackage{placeins}

\usepackage{textpos}
\usepackage{wrapfig}
\usepackage[spanish, es-tabla]{babel}

\usepackage{caption}
\setlength{\parindent}{0pt}
\hypersetup{hyperfootnotes=false}

\usepackage{url}
\usepackage{float}
\usepackage{tikz}
\usetikzlibrary{patterns}


\title[ICS-294]{Econometría ICS-294\\ \small{Clase 20: Raíces unitarias}}

\author{Francisco Alfaro Medina}
\institute[USM]
{Universidad Técnica Federico Santa María\\ % Your institution for the title page
	\medskip
	\textit{francisco.alfaro@usm.cl} % Your email address

}
\logo{

\includegraphics[width=3cm]{Figuras/logoUSM.png}
}
\date{}
\newcommand\Fontvi{\fontsize{9}{7.2}\selectfont}
\begin{document}


\frame{\titlepage}
\logo{}


\begin{frame}{Objetivos de la Clase}

Al finalizar esta clase, el estudiante será capaz de:

\begin{itemize}
    \item[\checkmark] Explicar qué es una \textbf{raíz unitaria} y por qué una caminata aleatoria es no estacionaria.
    \item[\checkmark] Argumentar por qué la regresión entre dos series $I(1)$ independientes produce un \textbf{$R^2$ alto y $t$ significativos falsos} (regresión espuria).
    \item[\checkmark] Derivar la regresión auxiliar del \textbf{test de Dickey--Fuller} ($\Delta y_t$ sobre $y_{t-1}$) y plantear $H_0:\rho=1$.
    \item[\checkmark] Elegir la especificación D-F adecuada (constante / constante $+$ tendencia) y comparar el estadístico con los \textbf{valores críticos D-F} (que no son los de la $t$ usual).
    \item[\checkmark] Aplicar el test ADF y resolver el problema diferenciando la serie.
\end{itemize}

\end{frame}


\begin{frame}{¿Qué es una Serie de Tiempo?}
Una \textbf{serie de tiempo} es un conjunto de observaciones cuantitativas ordenadas cronológicamente:
\[
Y_1, Y_2, ..., Y_T
\]
donde $Y_t$ es el valor observado en el instante $t$.

Las series de tiempo pueden presentar:
\begin{itemize}
    \item Tendencia
    \item Ciclos
    \item Estacionalidad
    \item Aleatoriedad
\end{itemize}
\end{frame}

\begin{frame}{Tipos de Modelos}
\begin{block}{Modelos deterministas}
Métodos simples de extrapolación sin considerar aleatoriedad. Ejemplo: promedio móvil.
\end{block}

\begin{block}{Modelos estocásticos}
Describen procesos aleatorios subyacentes usando variables aleatorias y distribuciones.
\end{block}
\end{frame}

% \begin{frame}{Estacionariedad}
% \textbf{Serie estacionaria:}
% \begin{itemize}
%     \item Media constante: $\mathbb{E}(Y_t) = \mathbb{E}(Y_{t+m})$
%     \item Varianza constante: $\text{Var}(Y_t) = \text{Var}(Y_{t+m})$
%     \item Covarianza depende solo del rezago: $\text{Cov}(Y_t, Y_{t+k}) = \text{Cov}(Y_{t+m}, Y_{t+m+k})$
% \end{itemize}

% \textbf{Serie no estacionaria:} sus características cambian con el tiempo. Puede hacerse estacionaria mediante diferenciación.
% \end{frame}

\begin{frame}{Componentes Clásicos de una Serie}
\begin{itemize}
    \item Tendencia ($T$): evolución a largo plazo.
    \item Ciclo ($C$): fluctuaciones de largo plazo no periódicas.
    \item Estacionalidad ($E$): variaciones periódicas de corto plazo.
    \item Aleatoriedad ($A$): componente residual, ruido.
\end{itemize}

\textbf{Modelos comunes:}
\begin{itemize}
    \item Aditivo: $Y(t) = T(t) + C(t) + E(t) + A(t)$
    \item Multiplicativo: $Y(t) = T(t) \cdot C(t) \cdot E(t) \cdot A(t)$
\end{itemize}
\end{frame}

\begin{frame}{Componentes}
\begin{figure}
    \centering
    \includegraphics[width=0.65\linewidth]{sem14//images/descomposicion_serie.png}
\end{figure}
\end{frame}

\begin{frame}
\frametitle{Ejemplo: Ventas Mensuales de Helados}

\begin{itemize}
    \item \textbf{Tendencia (T)}: crecimiento anual de las ventas por aumento poblacional.
    \item \textbf{Ciclo (C)}: caídas temporales por recesiones o eventos globales.
    \item \textbf{Estacionalidad (E)}: ventas más altas en verano, más bajas en invierno.
    \item \textbf{Aleatoriedad (A)}: variaciones impredecibles (semanas muy calurosas fuera de temporada, cortes eléctricos, etc.).
\end{itemize}

\begin{figure}
    \centering
    \includegraphics[width=0.75\linewidth]{sem14//images/venta_helado.png}
\end{figure}

\end{frame}

\begin{frame}
\frametitle{Proceso Estocástico}
\begin{itemize}
    \item Una serie de tiempo es una realización de un proceso estocástico: \(\{Y_t\}\)
    \item Un proceso estocástico es una colección de variables aleatorias indexadas por el tiempo.

    \[
\{ Y_t : t \in T \}
\]

donde \( T \) es el conjunto de tiempos (por ejemplo, \( T = \{1, 2, 3, \dots\} \) para datos discretos).

Cada \( Y_t \) es una variable aleatoria definida sobre un mismo espacio de probabilidad 
\((\Omega, \mathcal{F}, P)\), con una función de probabilidad conjunta que describe su comportamiento 
a lo largo del tiempo.

    \item Permite modelar la evolución temporal de fenómenos como precios, rendimientos, consumo, producción, etc.
\end{itemize}


\end{frame}

\begin{frame}{Proceso estocástico}
   \begin{itemize}
    \item \( \Omega \): representa el \textbf{espacio muestral}, es decir, el conjunto de todos los escenarios posibles del mundo 
    (todas las trayectorias que podrían ocurrir).

    \item \( \mathcal{F} \): es la \textbf{$\sigma$-álgebra} (sigma-álgebra), el conjunto de subconjuntos de \( \Omega \) sobre los cuales 
    se pueden definir eventos o sucesos medibles (por ejemplo: ``el PIB sube tres trimestres seguidos'').

    \item \( P \): es la \textbf{medida de probabilidad}, que asigna a cada evento \( A \in \mathcal{F} \) una probabilidad
    \( P(A) \in [0,1] \); en otras palabras, son las reglas del azar que gobiernan el proceso estocástico.
\end{itemize}

Así, cada variable \( Y_t \) es una función que asocia a cada escenario posible \( \omega \in \Omega \) 
un valor real \( Y_t(\omega) \), y la colección completa \( \{Y_t\} \) describe la evolución de esa variable a lo largo del tiempo.
 
\end{frame}

\begin{frame}{Ejemplo}
\small
\begin{itemize}
    \item Imaginemos que \( Y_t \) representa el \textbf{Producto Interno Bruto (PIB)} real trimestral de un país.
    
    \item El \emph{proceso estocástico} \( \{ Y_t : t \in T \} \) describe todas las posibles trayectorias 
    que el PIB podría seguir en distintos mundos posibles, es decir, bajo diferentes realizaciones 
    del azar \( \omega \in \Omega \).
    
    \item La \emph{serie de tiempo observada} corresponde a la trayectoria particular que efectivamente 
    ocurrió en nuestro mundo.
    
    \item En otras palabras, observamos una secuencia concreta de valores:
    \[
    y_t = Y_t(\omega_0), \quad t = 1, 2, \dots, T,
    \]
    para una realización específica \( \omega_0 \in \Omega \).
\end{itemize}
\end{frame}




\begin{frame}
\frametitle{Estacionariedad y Autocorrelación}
\begin{itemize}
   % \item Un proceso \(\{X_t\}\) es \textbf{estacionario en sentido estricto} si su distribución conjunta no cambia con el tiempo:
   % \[    (X_{t_1}, \ldots, X_{t_k}) \overset{d}{=} (X_{t_1+h}, \ldots, X_{t_k+h}) \quad \forall h \in \mathbb{Z} \]
    \item  Un proceso \(\{X_t\}\) es \textbf{estacionario } (débilmente o de segundo orden) si:
    \begin{itemize}
        \item Esperanza constante: \(\mathbb{E}[X_t] = \mu\)
        \item Varianza constante: \(\text{Var}(X_t) = \sigma^2 < \infty\)
        \item Covarianza depende solo del rezago: \(\gamma(h) = \text{Cov}(X_t, X_{t+h})\)
    \end{itemize}
    \item La \textbf{función de autocorrelación} es:
    \[
    \rho(h) = \frac{\gamma(h)}{\gamma(0)} = \frac{\text{Cov}(X_t, X_{t+h})}{\text{Var}(X_t)}
    \]
    \item En procesos estacionarios: \(\rho(h)\) depende solo del rezago \(h\), no del tiempo \(t\).
\end{itemize}
\end{frame}

\begin{frame}
\frametitle{Ruido Blanco}
\begin{itemize}
    \item Es un proceso estocástico \(\{\varepsilon_t\}\) con:
    \begin{itemize}
        \item Media cero: \(\mathbb{E}[\varepsilon_t] = 0\)
        \item Varianza constante: \(\text{Var}(\varepsilon_t) = \sigma^2\)
        \item No autocorrelación: \(\text{Cov}(\varepsilon_t, \varepsilon_{t+h}) = 0\) si \(h \neq 0\)
    \end{itemize}
    \item Es la base para construir modelos más complejos.
\end{itemize}
\end{frame}
\begin{frame}
\frametitle{Modelo Autorregresivo AR(p)}
\begin{itemize}
    \item Un modelo AR(p) se basa en los valores pasados de la serie:
    \[
    Y_t = \phi_1 Y_{t-1} + \cdots + \phi_p Y_{t-p} + \varepsilon_t \]

     Donde $\varepsilon_t$ es un ruido blanco.

    \item Ejemplo AR(1):
     \[
     Y_t = \phi Y_{t-1} + \varepsilon_t
     \]
\end{itemize}
 \end{frame}

\begin{frame}
\frametitle{Modelo de Media Móvil MA(q)}
\begin{itemize}
    \item El modelo MA(q) usa errores pasados:
    \[
    Y_t = \varepsilon_t + \theta_1 \varepsilon_{t-1} + \cdots + \theta_q \varepsilon_{t-q}
    \]
    \item Ejemplo MA(1):
    \[
    Y_t = \varepsilon_t + \theta \varepsilon_{t-1}
    \]
 \item Donde $\varepsilon_t$ es un ruido blanco.
\end{itemize}
\end{frame}


\begin{frame}{Correlograma }
    \begin{figure}
        \centering
        \includegraphics[width=0.85\linewidth]{sem14//images/correlograma1.png}
    \end{figure}
\end{frame}


\begin{frame}
\frametitle{Test de Breusch-Godfrey para correlación serial  de orden (1) {\bf del error}}

\begin{itemize}
  \item Para series de tiempo existen varios test.
  \item Con la idea de una autocorrelación de orden 1:
  \[
  y_t = x_t \beta + \epsilon_t
  \]
  \[
  \epsilon_t = \rho_1 \epsilon_{t-1}+v_t
  \]
  \item Test LM para testear \( H_0: \rho_1 = 0 \) (si hay un solo rezago, también se puede hacer test t).
  \item Bajo \( H_0 \) cierta (no hay autocorrelación de orden 1) tenemos:
 % \begin{itemize}
    % \item Si \( \text{LM}_{\text{empírico}} > \chi^2_{\text{tabla}} \) \(\rightarrow\) ¡Rechazo \( H_0 \)!
  %\end{itemize}
\end{itemize}

\end{frame}
\begin{frame}
\frametitle{Pasos Breusch-Godfrey}

\begin{itemize}
  \item La prueba de hipótesis la escribimos:
  \[
  H_0 : \rho_1 = 0
  \]
  \[
  H_1 : \rho_1 \neq 0
  \]
  \item Pasos a seguir:
  \begin{enumerate}
    \item Estimar el modelo original por MCO y obtener los residuos \(\hat{\epsilon}\).
    \item Estimar una regresión auxiliar de \(\hat{\epsilon}_t\) sobre el rezago: \(\hat{\epsilon}_t = \rho_0 + \rho_1 \hat{\epsilon}_{t-1}+ \nu_t\) .
    \item Quedarse con el \(R^2\) y construir el estadístico LM de prueba:
    \[
    \text{LM} = nR^2_\epsilon
    \]
    \item Se compara con una \(\chi^2_p\) donde \(p\) es el número de rezagos (1 en este caso).
    \item Si \(\text{LM} > \chi^2_p\) al \(\alpha \%\) \(\rightarrow\) rechazo la hipótesis nula de ausencia de autocorrelación. ¡Tenemos problema!
  \end{enumerate}
\end{itemize}

\end{frame}

\begin{frame}
\frametitle{Estimación por MCO con errores robustos a la autocorrelación}

\begin{itemize}
  \item Al igual que con la heteroscedasticidad, podemos estimar por MCO y obtener errores robustos a la autocorrelación.
  \item \textbf{Newey} y \textbf{West} proponen una forma de estimar los errores robusta a la autocorrelación y la heteroscedasticidad.
  \item Llamamos a estos errores \textbf{HAC} por sus siglas en inglés.
\end{itemize}

\end{frame}

\begin{frame}
\frametitle{Errores Robustos de \textbf{Newey} y \textbf{West}}

\begin{itemize}
    \item La metodología de Newey y West permite obtener errores estándar robustos en presencia de:
    \begin{itemize}
        \item \textbf{Heterocedasticidad} (varianza no constante)
        \item \textbf{Autocorrelación} (errores correlacionados en el tiempo de orden q)
    \end{itemize}

    \item Es una extensión del estimador “sandwich”:
    \[
    \widehat{\text{Var}}(\hat{\beta}) = (X'X)^{-1} \hat{\Omega} (X'X)^{-1}
    \]
          \item  \( \hat{\Omega} \) es una estimación robusta de la matriz de varianzas-covarianzas de los errores, para obtenerla se estima el modelo por MCO y se obtienen los residuos \( \hat{\epsilon}_t\).
 

\item Considera pesos aplicados a cada rezago, que disminuyen linealmente. Estos pesos reducen la importancia de correlaciones lejanas, haciendo que el estimador sea más estable.
   
\end{itemize}
\end{frame}

\begin{frame}
\frametitle{Series de Tiempo con Tendencia Determinística}
\begin{itemize}
    \item Una serie tiene \textbf{tendencia determinística} cuando su valor esperado cambia sistemáticamente con el tiempo:
    \[
    Y_t = \beta_0 + \beta_1 t + \varepsilon_t
    \]
    \item Características:
    \begin{itemize}
        \item La media \(\mathbb{E}[Y_t] = \beta_0 + \beta_1 t\) crece o decrece linealmente.
        \item El componente \(\varepsilon_t\) suele ser ruido blanco o estacionario.
        \item La tendencia es predecible (determinada por el tiempo).
    \end{itemize}
    \item Ejemplos típicos:
    \begin{itemize}
        \item Ventas acumuladas con crecimiento constante.
        \item Temperatura media anual en contextos de cambio climático.
    \end{itemize}
   
\end{itemize}
\end{frame}


\begin{frame}
\frametitle{Inclusión de tendencias en los modelos de series temporales}

\begin{itemize}
  \item La "tendencia" de las series es como una variable omitida!
  \item Una manera de solucionar esto es incluir en el modelo una variable que recoja el crecimiento temporal: la tendencia.
  \item Una tendencia lineal \( t \) es una variable que crecerá una unidad cada año, lo más común empezando en 1.
\begin{center}
\begin{tabular}{c c}
\hline
Año & t \\
\hline
1940 & 1 \\
1941 & 2 \\
1942 & 3 \\
\vdots & \vdots \\
1960 & 20 \\
\hline
\end{tabular}
\end{center}
  \item Se pueden ajustar tendencias cuadráticas (\( t^2 \)), o si la variable dependiente está en logs, la tendencia lineal será interpretada como tendencia exponencial.
\end{itemize}

\end{frame}
\begin{frame}
\frametitle{Inclusión de tendencias en los modelos de series temporales}

\begin{itemize}
  \item Cuando incluimos una tendencia determinística, el modelo que estimaremos será:
  \[
  y_t = \alpha_0 + \alpha_1 x_t + \alpha_2 t + \epsilon_t
  \]
  \item \(\alpha_2\) nos dice cuánto varía anualmente \( y_t \):
  \begin{itemize}
    \item \(\alpha_2 > 0\) y significativo: tendencia creciente.
    \item \(\alpha_2 < 0\) y significativo: tendencia decreciente.
  \end{itemize}
  \item Ahora la ecuación la debemos interpretar como “libre de tendencia”.
  \item \(\alpha_1\) nos dirá cuál es la relación entre \( x_t \) y \( y_t \) cuando les “quitamos” la tendencia.
\end{itemize}

\end{frame}
\begin{frame}
\frametitle{Persistencia en modelos de series de tiempo}

\begin{itemize}
  \item Muchas veces, las variables del modelo tienen persistencia: $ y_t = \rho y_{t-1} + \epsilon_t $.
  \item Si \(\rho = 1\), es una \textbf{caminata aleatoria}:
  \begin{itemize}
    \item Tiene una raíz unitaria o que es \textbf{integrada de orden 1.}
       \item Al igual que la tendencia 'determinística' que vimos antes, cuando tenemos una variable que se comporta de esta manera, vamos a tener que ajustarlo. Pero no va a alcanzar con agregar tendencia en este caso.
  \end{itemize}
  \item Cuando \(|\rho| < 1\)
  \begin{itemize}
    \item No hay tendencia estocástica.
    \item En este caso, podemos usar MCO sin inconveniente.
  \end{itemize}
\end{itemize}

\end{frame}

\begin{frame}{Correlograma caminata aleatoria }


    \begin{figure}
        \centering
        \includegraphics[width=0.85\linewidth]{sem14//images/correlograma2.png}
    \end{figure}
\end{frame}
\begin{frame}
\frametitle{Series altamente persistentes}

\begin{itemize}
  \item Si \(\rho = 1\), usar estas variables va a introducir correlación espuria entre las series.
  \item Los datos tienen una tendencia en el tiempo, pero no es determinística (crece de forma constante) sino que es estocástica (crece pero de forma variante).
  \item Lo vamos a solucionar tomando la primera diferencia de la serie problemática:
  $
  \Delta y_t = y_t - y_{t-1} = \epsilon_t 
  $ y se corre la regresión con la diferencia y no el nivel.
  \item Ejemplo: 
$
  salario_t = \beta_0 + \beta_1 productividad_t + \epsilon_t 
$
  Supongan que tanto salario como productividad siguen un proceso de raíz unitaria (caminata aleatoria).
  \item Para poder estimar su relación, vamos a transformar el modelo a primeras diferencias.
  Con esto, nos sacamos de encima la persistencia.
\end{itemize}
$
  \Delta salario_t = \beta_0 + \beta_1 \Delta productividad_t + u_t
$
\end{frame}
\begin{frame}
\frametitle{Exploración gráfica - I(1)}

\begin{itemize}
  \item AR(1) con \(\rho = 1\)

             \begin{figure}
        \centering
        \includegraphics[width=0.7\linewidth]{sem14/images/ar1.png}
    \end{figure}


  \item La persistencia es muy alta
\end{itemize}

\end{frame}

\begin{frame}
\frametitle{Exploración gráfica - Primera diferencia de una I(1)}

\begin{itemize}
  \item La serie con orden de integración 1 la diferenciamos, y ¿qué obtenemos?

             \begin{figure}
        \centering
        \includegraphics[width=0.7\linewidth]{sem14/images/i1.png}
    \end{figure}
  
  \item Muchas series macroeconómicas siguen este tipo de procesos
\end{itemize}

\end{frame}

\begin{frame}{Raíces unitarias, tendencia estocástica y regresión espuria}
    \begin{itemize}
        \item Queremos estimar el modelo:
        \[
        Y_t = \beta_0 + \beta_1 X_t + u_t
        \]
        \item Supongamos que el proceso de cada variable es el siguiente:
        \[
        Y_t = \alpha_1 + \rho_1 Y_{t-1} + \epsilon_t
        \]
        \[
        X_t = \alpha_2 + \rho_2 X_{t-1} + v_t
        \]
        \item Si $-1 < \rho < 1$, entonces los procesos son AR(1) estacionarios y se puede hacer la regresión sin problema. 
        \item Si $\rho = 1$ la serie es una caminata aleatoria.
     Esto hace que se encuentren relaciones espurias entre $x$ e $y$: $\beta_1$ significativo cuando en realidad no hay relación.
    \end{itemize}
\end{frame}

\begin{frame}[allowframebreaks]{¿Por qué dos $I(1)$ independientes dan regresión espuria?}

\begin{alertblock}{El fenómeno de la regresión espuria (Granger--Newbold, 1974)}
Si $y_t$ y $x_t$ son caminatas aleatorias \textbf{independientes} (no hay relación verdadera, $\beta_1=0$) y aun así regresamos $y_t=\beta_0+\beta_1 x_t+u_t$:
\begin{itemize}
  \item $\hat\beta_1$ tiende a parecer ``significativo'' con altísima frecuencia.
  \item El $R^2$ es alto.
  \item El estadístico $t$ \textbf{no} converge: crece con $\sqrt{T}$.
\end{itemize}
\end{alertblock}

\begin{exampleblock}{Intuición de la demostración}
\textbf{Paso 1.} Cada serie acumula sus choques: $y_t=\sum_{j\le t}\epsilon_j$, $x_t=\sum_{j\le t}v_j$, con $\epsilon\perp v$. Ambas ``vagan'' lejos de su media sin volver.

\textbf{Paso 2.} El residuo del modelo, $u_t=y_t-\beta_0-\beta_1 x_t$, \emph{también} es $I(1)$: hereda la no estacionariedad. Por tanto $u_t$ está fuertemente autocorrelacionado y \textbf{no} es ruido blanco.

\framebreak

\textbf{Paso 3.} El cálculo del error estándar de MCO supone $u_t$ con varianza constante y sin autocorrelación. Como aquí $\text{Var}(u_t)$ crece con $t$, ese error estándar está \textbf{subestimado}:
\[
t = \frac{\hat\beta_1}{\widehat{\text{se}}(\hat\beta_1)} \xrightarrow[\text{} ]{}\ \text{diverge (no tiene distribución } t).
\]
\textbf{Paso 4.} En consecuencia rechazamos $H_0:\beta_1=0$ casi siempre, ¡aunque sea cierta! El $R^2$ alto solo refleja que dos series que ``derivan'' parecen moverse juntas por azar.

\textbf{Moraleja:} con series $I(1)$, un $t$ grande y un $R^2$ alto \textbf{no} son evidencia de relación. Hay que verificar el orden de integración \emph{antes} de regresar. \hfill $\square$
\end{exampleblock}

\end{frame}

\begin{frame}{Síntoma práctico de regresión espuria}

\begin{block}{Regla de Granger--Newbold}
Sospeche de una regresión espuria cuando:
\[
R^2 > d_{\text{Durbin--Watson}}.
\]
Es decir, cuando el ajuste aparente es alto pero el estadístico Durbin--Watson es muy bajo (residuos fuertemente autocorrelacionados, típico de $u_t$ que es $I(1)$).
\end{block}

\begin{itemize}
  \item La solución es la misma de la Clase 19: \textbf{diferenciar} las series $I(1)$ y regresar en diferencias, donde $\Delta y_t$ y $\Delta x_t$ son $I(0)$.
\end{itemize}

\end{frame}

\begin{frame}{Solución al problema de regresión espuria por raíz unitaria: diferenciar las variables}
    \begin{itemize}
        \item Si $\rho = 1$, en lugar de estimar el modelo original, hacemos lo siguiente:
        \begin{enumerate}
            \item Diferenciamos las variables:
            \[
            \Delta Y_t = Y_t - Y_{t-1}, \quad \Delta X_t = X_t - X_{t-1}
            \]
            \item Hacemos la regresión en diferencias:
            \[
            \Delta Y_t = \beta_0 + \beta_1 \Delta X_t + u_t
            \]
        \end{enumerate}
        \item Intuición:
        \[
        Y_t = \alpha_1 + Y_{t-1} + \epsilon_t \implies \text{altamente persistente}
        \]
        \[
        \Delta Y_t = \alpha_1 + \epsilon_t \implies \text{¡estacionario!}
        \]
       % \item Por lo que las series en diferencias ya no tienen el problema.
    \end{itemize}
\end{frame}
\begin{frame}{¿Cómo saber si $\rho = 1$? Test visual y Test de D-F}
    \begin{itemize}
        \item Naturalmente uno querría estimar el siguiente modelo para saber si la serie tiene una raíz unitaria:
        \[
        y_t = \alpha + \rho y_{t-1} + \epsilon_t
        \]
        \item Y hacer una prueba de hipótesis sobre $\rho = 1$.
        \item El problema es que si $\rho = 1$ la estimación de $\rho$ no es válida y el estadístico t de contraste no seguirá una distribución paramétrica conocida, ni siquiera en muestras grandes.
        \item En otras palabras, estaríamos usando para hacer la prueba un estadístico que no podemos estimar consistentemente y por eso la prueba no es válida.
    \end{itemize}
\end{frame}
\begin{frame}{Test de raíces unitarias}
    \begin{itemize}
        \item Dickey-Fuller (D-F) derivan distribuciones asintóticas válidas para estos casos que nos van a permitir hacer inferencia válida sobre $\rho$:
        \begin{eqnarray}
 y_t = \alpha + \rho y_{t-1} + \epsilon_t\\
 y_t - y_{t-1} = \alpha + \rho y_{t-1} - y_{t-1} + \epsilon_t\\
  \Delta y_t = \alpha + (\rho - 1) y_{t-1} + \epsilon_t\\
   \Delta y_t = \alpha + \phi y_{t-1} + \epsilon_t
        \end{eqnarray}

   
        \item Para esta especificación tenemos las tablas de distribución de D-F y por tanto podemos testear:
        \begin{itemize}
            \item $H_0: \phi = 0 \implies \rho = 1 \implies \text{Raíz unitaria (1)}$
            \item $H_1: \phi < 0 \implies \rho < 1 \implies \text{Serie estacionaria (2)}$
        \end{itemize}
        \item Construimos el estadístico t para la significación de $\phi$ y luego comparamos con los valores críticos de las tablas D-F.
    \end{itemize}
\end{frame}

\begin{frame}{Idea clave: la reparametrización de Dickey--Fuller}

\begin{alertblock}{Por qué se reparametriza}
No podemos testear $\rho=1$ con la $t$ usual sobre $y_t=\alpha+\rho y_{t-1}+\epsilon_t$, porque bajo $\rho=1$ el regresor $y_{t-1}$ es $I(1)$ y el estadístico \textbf{no} sigue una $t$ de Student ni siquiera asintóticamente.
\end{alertblock}

\begin{exampleblock}{Demostración de la equivalencia $\rho=1 \Leftrightarrow \phi=0$}
\textbf{Paso 1.} Restamos $y_{t-1}$ a ambos lados:
\[
y_t - y_{t-1} = \alpha + \rho y_{t-1} - y_{t-1} + \epsilon_t.
\]
\textbf{Paso 2.} Factorizamos $y_{t-1}$ y definimos $\phi \equiv \rho-1$:
\[
\Delta y_t = \alpha + (\rho-1)\,y_{t-1} + \epsilon_t = \alpha + \phi\, y_{t-1} + \epsilon_t.
\]
\textbf{Paso 3.} Por tanto:
\[
\rho=1 \iff \phi=0 \quad(\text{raíz unitaria}), \qquad \rho<1 \iff \phi<0 \quad(\text{estacionaria}).
\]
La ventaja: ahora $\Delta y_t$ es $I(0)$ y D-F tabularon la distribución del estadístico $\tau=\hat\phi/\widehat{\text{se}}(\hat\phi)$ bajo $H_0$ mediante simulación. Es un test \textbf{de una cola por la izquierda}: rechazamos si $\tau$ es \emph{muy negativo}. \hfill $\square$
\end{exampleblock}

\end{frame}

\begin{frame}{Test Dickey--Fuller Aumentado (ADF)}

\begin{itemize}
  \item El test D-F simple supone que $\epsilon_t$ es ruido blanco. Si los errores tienen \textbf{autocorrelación}, se aumenta la regresión con rezagos de $\Delta y_t$:
  \[
  \Delta y_t = \alpha + \phi\, y_{t-1} + \delta t + \sum_{i=1}^{p} \gamma_i\, \Delta y_{t-i} + \epsilon_t
  \]
  \item Los rezagos $\Delta y_{t-i}$ ``limpian'' la autocorrelación para que $\epsilon_t$ vuelva a ser ruido blanco; la hipótesis y los valores críticos son los \emph{mismos} que en D-F.
  \item $H_0:\phi=0$ (raíz unitaria, $I(1)$) vs.\ $H_1:\phi<0$ (estacionaria).
  \item El número de rezagos $p$ se elige con criterios de información (AIC/BIC) o reglas prácticas.
\end{itemize}

\end{frame}

\begin{frame}{Distribuciones de D-F}
    \begin{itemize}
        \item No existe una única distribución de D-F: dependerá de contra qué estemos comparando la hipótesis nula de existencia de raíz unitaria.
        \item Esto básicamente depende de cómo se comportan los datos y por tanto de cómo esté especificado el modelo.
        \item En otras palabras, en cada caso debemos pensar que nuestra hipótesis nula es que la serie tiene raíz unitaria y la alternativa puede ser:
        \begin{enumerate}
            \item Que sea estacionaria pero se mueva en el tiempo en torno a una media distinta de cero.
            \item Que sea estacionaria pero tenga una tendencia lineal en el tiempo (creciente o decreciente).
        \end{enumerate}
    \end{itemize}
\end{frame}
\begin{frame}{Distribuciones de D-F}
    \begin{itemize}
        \item D-F distinguen entre dos casos:
        \begin{enumerate}
            \item Solo constante: bajo la alternativa la serie no tiene raíz unitaria pero tiene una media distinta de cero
            \[
            y_t = \alpha + \rho y_{t-1} + \epsilon_t
            \]
            \item Constante y tendencia lineal: bajo la hipótesis alternativa la serie tiene tendencia pero es lineal en el tiempo
            \[
            y_t = \alpha + \rho y_{t-1} + \delta t + \epsilon_t
            \]
        \end{enumerate}
        \item En cada caso hay una tabla distinta con valores críticos distintos.
    \end{itemize}
\end{frame}
\begin{frame}{Valores críticos de las distribuciones de D-F}
    \begin{itemize}
        \item Los valores críticos varían según la hipótesis alternativa contra la que se quiere testear y según el tamaño de la muestra.

             \begin{figure}
        \centering
        \includegraphics[width=0.7\linewidth]{sem14/tabla.png}
    \end{figure}
        
        \item Para ver cuál es el modelo alternativo que tiene sentido, hay que empezar por mirar los datos.
    \end{itemize}
\end{frame}

\begin{frame}{Valores cr\'iticos de las distribuciones de D-F}
    Los valores críticos para la prueba de Dickey-Fuller suelen obtenerse a través de simulaciones Monte Carlo, reflejando mejor la distribución real del estadístico bajo la hipótesis nula. Esto implica que los valores críticos se derivan empíricamente y son específicos a la estructura del modelo de la prueba.
\end{frame}

%\begin{frame}{Exploración gráfica - AR(1)}
%    \begin{itemize}
%        \item AR(1) con $\rho = 0.5$
%    \end{itemize}
%             \begin{figure}
%        \centering
%        \includegraphics[width=0.7\linewidth]{sem14/images/ar1.png}
%    \end{figure}
%\end{frame}

\begin{frame}
\frametitle{El PIB Chileno}

\begin{itemize}
    \item Queremos saber cu\'al es el orden de integraci\'on del PIB chileno.
    \item Comenzamos por una exploraci\'on gr\'afica

                 \begin{figure}
        \centering
        \includegraphics[width=0.6\linewidth]{sem14/images/azul.png}
    \end{figure}
    
   
\end{itemize}

\end{frame}
\begin{frame}
\frametitle{El PIB Chileno - Primera Diferencia}

\begin{itemize}
    \item Primera diferencia del PIB

                     \begin{figure}
        \centering
        \includegraphics[width=0.6\linewidth]{sem14/images/rojo.png}
    \end{figure}
    
    \item ¿Qu\'e impresi\'on nos da?
\end{itemize}

\end{frame}
\begin{frame}
\frametitle{El PIB Chileno - Ambas Gráficas}

\begin{itemize}
    \item Poniendo ambas gr\'aficas juntas
\end{itemize}

                 \begin{figure}
        \centering
        \includegraphics[width=0.6\linewidth]{sem14/images/azulroj.png}
    \end{figure}

\end{frame}
\begin{frame}
\frametitle{El PIB Chileno - Test de DF}

\begin{itemize}
    \item Realizamos el test de D-F para ver si se trata de una serie integrada de orden 1
    \item Dada nuestra exploraci\'on gr\'afica, corresponder\'ia usar como modelo alternativo el que incluye una tendencia temporal
$$\Delta PIB_t = \alpha + \phi PIB_{t-1} + \delta t +\epsilon_t$$
                     
    \begin{figure}
        \centering
        \includegraphics[width=0.7\linewidth]{sem14/images/pib.png}
    \end{figure}
    
    \item ¿Con qu\'e valor cr\'itico comparamos el estad\'istico t?
\end{itemize}

\end{frame}
\begin{frame}
\frametitle{El PIB Chileno - Test de DF}

\begin{itemize}
    \item Vamos a usar los valores de tabla de DF
    \item $\hat{\phi} = -0.0373029$ y $\hat{\sigma_{\phi}}=0.0228403$ 
    \item Como nuestra muestra es de 60, podemos usar los valores cr\'iticos de la de 50:
    \begin{itemize}
        \item VC al 5\%: -3.50
        \item $| -1.63 | < | -3.50 |$
        \item No rechazo $H_0$
    \end{itemize}
    \item Evidencia a favor de que el PIB chileno tiene una ra\'iz unitaria: no es una serie estacionaria
\end{itemize}

\end{frame}

%%%%%%%%%%%%%%%%%%%%%%%%%% EJERCICIOS RESUELTOS %%%%%%%%%%%%%%%%%%%%%%%%%%
\section{Ejercicios}

\begin{frame}{Ejercicio Resuelto 1}

\begin{block}{Ejercicio 1 --- Test D-F sobre el PIB real de Chile (adaptado Certamen 3)}
Con $T=25$ observaciones anuales del PIB real, y dada una exploración gráfica que sugiere tendencia, se estima la regresión D-F con constante y tendencia:
\[
\Delta y_t = \alpha + \phi\, y_{t-1} + \delta t + \epsilon_t.
\]
El estadístico para $\phi$ resulta $\tau = -2.93$. El valor crítico D-F al 5\% (con tendencia, $T\approx 25$) es $-3.60$. ¿Tiene el PIB raíz unitaria?
\end{block}

\pause

\begin{exampleblock}{Solución}
\textbf{Paso 1 (hipótesis):} $H_0:\phi=0$ (raíz unitaria, $I(1)$) vs.\ $H_1:\phi<0$ (estacionaria con tendencia).

\textbf{Paso 2 (regla):} test de una cola por la izquierda. Rechazo $H_0$ solo si $\tau < \text{VC}$, es decir si $\tau$ es \emph{más negativo} que $-3.60$.

\textbf{Paso 3 (comparación):} $\tau=-2.93 > -3.60$ $\Rightarrow$ \textbf{no se rechaza} $H_0$ (equivalente: $|-2.93|<|-3.60|$).

\textbf{Paso 4 (conclusión):} hay evidencia de \textbf{raíz unitaria}: el PIB real es $I(1)$, no estacionario. Para usarlo en regresión, habrá que diferenciarlo.
\end{exampleblock}

\end{frame}

\begin{frame}{Ejercicio Resuelto 2}

\begin{block}{Ejercicio 2 --- Dos series con distinto orden de integración (adaptado Certamen 3)}
Se quiere explicar el PIB per cápita ($\text{PIBpc}_t$) con la inversión neta ($\text{InvNeta}_t$), datos 1972--2016. Los tests D-F con tendencia entregan: para $\text{PIBpc}_t$, $\phi$ con $p=0.139$; para $\text{InvNeta}_t$, $\phi<0$ con $p=0.039$. ¿Qué orden de integración tiene cada una y cómo debe especificarse la regresión para evitar sesgos?
\end{block}

\pause

\begin{exampleblock}{Solución}
\textbf{Paso 1:} $\text{PIBpc}_t$: $p=0.139>0.05$ $\Rightarrow$ no rechazo $H_0$ $\Rightarrow$ \textbf{raíz unitaria}, $I(1)$.

\textbf{Paso 2:} $\text{InvNeta}_t$: $p=0.039<0.05$ $\Rightarrow$ rechazo $H_0$ $\Rightarrow$ \textbf{estacionaria alrededor de tendencia}, $I(0)$ con tendencia.

\textbf{Paso 3:} regresar el nivel $I(1)$ contra un $I(0)$ produciría regresión espuria. Hay que igualar órdenes: diferenciar la serie $I(1)$ e incluir la tendencia determinística:
\[
\Delta \text{PIBpc}_t = \beta_0 + \beta_1\, \text{InvNeta}_t + \beta_2\, t + u_t,
\]
donde $\Delta \text{PIBpc}_t$ es $I(0)$ e $\text{InvNeta}_t$ es $I(0)$ (en torno a tendencia). \hfill $\square$
\end{exampleblock}

\end{frame}

\begin{frame}[fragile]{Ejemplo Aplicado en Python --- Test ADF del IMACEC}

\begin{exampleblock}{Dickey--Fuller Aumentado con statsmodels}
\begin{verbatim}
import pandas as pd
from statsmodels.tsa.stattools import adfuller

datos = pd.read_csv("imacec.csv")
serie = datos["imacec"]

# ADF en nivel: regtype 'ct' = constante + tendencia
adf_nivel = adfuller(serie, regression="ct", autolag="AIC")
print("ADF nivel:  estadístico =", adf_nivel[0],
      " p-valor =", adf_nivel[1])

# ADF en primera diferencia
adf_diff = adfuller(serie.diff().dropna(), regression="c")
print("ADF Delta:  estadístico =", adf_diff[0],
      " p-valor =", adf_diff[1])
# Tipico: no rechaza en nivel (I(1)),
#         rechaza en diferencia (I(0))
\end{verbatim}
\end{exampleblock}

\end{frame}

\begin{frame}{Ejercicios Propuestos}

\begin{block}{Ejercicio 3 (Propuesto --- rutinario)}
Un test D-F \emph{solo con constante} sobre la inflación mensual del IPC ($T=120$) entrega $\tau=-4.10$. El valor crítico al 5\% es $-2.89$. ¿Es la inflación estacionaria? Plantee $H_0$, $H_1$ y concluya.
\end{block}

\begin{block}{Ejercicio 4 (Propuesto --- intermedio)}
A partir de $y_t=\alpha+\rho y_{t-1}+\delta t+\epsilon_t$, derive paso a paso la forma D-F $\Delta y_t=\alpha+\phi y_{t-1}+\delta t+\epsilon_t$ e indique qué valor de $\phi$ corresponde a la hipótesis nula de raíz unitaria.
\end{block}

\begin{block}{Ejercicio 5 (Propuesto --- desafiante)}
Dos analistas regresan el dólar observado sobre el precio del cobre (ambas $I(1)$) y obtienen $R^2=0.88$ con $t=9.5$, pero Durbin--Watson $=0.21$. Discuta si la relación es creíble, qué fenómeno está ocurriendo y cómo reespecificaría el modelo.
\end{block}

\end{frame}

\begin{frame}{Resumen}

\begin{block}{Puntos Clave de la Clase 20}
\begin{enumerate}
    \item Una serie con \textbf{raíz unitaria} ($\rho=1$) es $I(1)$, no estacionaria, con tendencia estocástica.
    \item Regresar dos series $I(1)$ independientes produce \textbf{regresión espuria}: $R^2$ alto y $t$ ``significativos'' falsos, porque el residuo es $I(1)$ y el error estándar está subestimado.
    \item El \textbf{test D-F} reparametriza a $\Delta y_t=\alpha+\phi y_{t-1}(+\delta t)+\epsilon_t$ y contrasta $H_0:\phi=0$ ($\rho=1$). Es de una cola por la izquierda.
    \item Los \textbf{valores críticos} son los de D-F (más negativos que los de la $t$), y dependen de la especificación (con/sin tendencia) y de $T$.
    \item Si hay autocorrelación se usa \textbf{ADF}; si se detecta $I(1)$, se \textbf{diferencia} y se estima en diferencias.
\end{enumerate}
\end{block}

\begin{alertblock}{Para la próxima clase}
Revisar: ¿qué pasa cuando los \emph{errores} están correlacionados en el tiempo? $\Rightarrow$ Autocorrelación serial, tests de Breusch--Godfrey y Durbin--Watson, y errores HAC (Clase 21).
\end{alertblock}

\end{frame}

\begin{frame}[shrink=10]
\frametitle{Conclusiones y Repaso}

\begin{itemize}
    \item Cuando trabajamos con series de tiempo debemos comenzar por ver si las series tienen tendencias temporales: determin\'isticas o estoc\'asticas
    \item Si es determinista, agregamos una tendencia:
    \[
    Y_t = \beta_0 + \beta_1 X_t + \alpha t + u_t
    \]
    \item Interpretaci\'on de $\beta_1$ ahora?
    \item Para saber si tiene una tendencia estoc\'astica (ra\'iz unitaria) procedemos a realizar el test de Dickey-Fuller
    \item Este test tiene valores cr\'iticos espec\'ificos y los mismos dependen del modelo alternativo que nos estemos planteando (con o sin tendencia determin\'istica)
    \item Para determinar el modelo alternativo comenzamos por una exploraci\'on gr\'afica
    \item Si encontramos que la serie tiene una ra\'iz unitaria (no rechazamos $H_0$) vamos a tener que diferenciar la serie y estimar nuestro modelo en diferencias.
\end{itemize}

\end{frame}

% \begin{frame}
% \frametitle{Conclusión}

% \begin{itemize}
%   \item Las series de tiempo muchas veces presentan tendencias:
%   \begin{itemize}
%     \item Determinísticas: cuando la serie varía de forma constante (no aleatoria) con el tiempo.
%     \item Estocásticas: cuando la serie varía de forma variable (aleatoria) con el tiempo.
%   \end{itemize}
%   \item Siempre que haya tendencias, las estimaciones por MCO serán sesgadas: correlaciones espurias (variables omitidas).
%   \item Si la tendencia es determinística: incluimos un regresor que varíe constantemente con el tiempo, por ejemplo, una tendencia lineal.
%   \item Si la tendencia es estocástica: tenemos que usar la primera diferencia de la serie para estimar.
%   \item ¿Cómo sabemos si la serie es integrada de orden 1?
%   \begin{itemize}
%     \newline Se utiliza el test de Dickey-Fuller.
%   \end{itemize}
% \end{itemize}

% \end{frame}


% \begin{frame}
% \frametitle{Procesos Integrados}
% \begin{itemize}
%     \item Muchas series no son estacionarias.
%     \item Una serie es integrada de orden \(d\) si se vuelve estacionaria al diferenciarla \(d\) veces.
%     \item Diferencia de primer orden:
%     \[
%     \nabla Y_t = Y_t - Y_{t-1}
%     \]
%     \item Si \(\nabla^d Y_t\) es estacionaria, entonces \(Y_t \sim I(d)\).
% \end{itemize}
% \end{frame}

% \begin{frame}
% \frametitle{Modelo ARIMA(p,d,q)}
% \begin{itemize}
%     \item Combina:
%     \begin{itemize}
%         \item AR(p): Autorregresivo
%         \item I(d): Diferencias para estacionariedad
%         \item MA(q): Media móvil
%     \end{itemize}
%     \item Notación:
%     \[
%     \phi(L)(1 - L)^d Y_t = \theta(L) \varepsilon_t
%     \]
%     \item \(L\) es el operador de rezago: \(L Y_t = Y_{t-1}\)
% \end{itemize}
% \end{frame}

\end{document}
